\clearpage
\section{Vector Spaces}

\subsection{$\mathbb{R}^n$ and $\mathbb{C}^n$}

These exercises are really easy but I am always scared that my proofs of these are not rigorous enough....

I do odds.

\begin{enumerate}
    \item[1] Show that $𝛼 + 𝛽 = 𝛽 + 𝛼$ for all $𝛼, 𝛽 ∈ \mathbb{C}$.
    
    \begin{proof}
        Let $\alpha = a + bi$ and $\beta = c + di$ where $a, b, c, d \in \mathbb{R}$.

        Then we have 
        \begin{align*}
            \alpha + \beta &= \\
            a + bi + c + di &= \\ 
            (a + c) + (b + d)i &= \\
        \end{align*}
        Because addition is commutative under the reals, we can swap the terms.
        \begin{align*}
            (c + a) + (d + b)i &= \\
            c + di + a + bi &= \\
            \beta + \alpha  
        \end{align*}
    \end{proof}

    \item[3] Show that $(𝛼𝛽)𝜆 = 𝛼(𝛽𝜆)$ for all $𝛼, 𝛽, 𝜆 ∈ \mathbb{C}$.

    Distribute, multiply, then rearrange terms using the communitivity of the reals to get the answer. It is quite tedious, so I won't do it here.

    \item[5] Show that for every $𝛼 ∈ \mathbb{C}$, there exists a unique $𝛽 ∈ \mathbb{C}$ such that $𝛼 + 𝛽 = 0$.
    
    \begin{proof}
        Let $\alpha = a + bi$. Then, let $\beta = -a - bi$. We can observe that 
        \begin{align*}
            \alpha + \beta &= \\
            a + bi - a - bi &= \\
            (a - a) + i(b - b) &= 0
        \end{align*}
    \end{proof}

    \item[7] Show that $\frac{-1 + \sqrt{3}i}{2}$ is a cube root of 1 (meaning that its cube equals 1)
    
    Again, this is a tedious problem where you just multiply it out. Of course, I can represent this number as $e^{\frac{\pi}{3}i}$ but that's cheating.

    \item[9] Find $𝑥 ∈ 𝐑^4$ such that

    $(4, −3, 1, 7) + 2𝑥 = (5, 9, −6, 8)$

    \begin{proof}
        This is vector addition, so we can apply the definition. 

        \begin{align*}
            (4, -3, 1, 7) + 2x &= (5, 9, -6, 8) \\
            (4 + 2x_1, -3 + 2x_2, 1 + 2x_3, 7 + 2x_4) &= (5, 0, -6, 8)
        \end{align*}
        For each term in the list, the number must hold, so in essence we get four different equations:
        \begin{align*}
            4 + 2x_1 &= 5 \\
            -3 + 2x_2 &= 0 \\
            1 + 2x_3 &= -6 \\
            7 + 2x_4 &= 8 \\
        \end{align*}
        Each of these linear equations yields only one solution, so the only solution for the equation is 
        \boxed{x = (\frac{1}{2}, -\frac{3}{2}, -\frac{7}{2}, \frac{1}{2})}
    \end{proof}
\end{enumerate}

\subsection{Definition of a Vector Space}

This one introduces the concept of a vector space in an abstract manner and gives a few unorthodox examples.

\begin{enumerate}
    \item[1] Prove that $−(−𝑣) = 𝑣$ for every $𝑣 ∈ V$

    \begin{proof}
        $-(-v)$ is the additive inverse of $-v$. Therefore,
        
        \begin{align*}
            -(-v) + (-v) &= 0 \\
            -(-v) + (-v + v) &= v \\
            -(-v) + 0 &= v \\
            -(-v) = v 
        \end{align*}
        As desired.
    \end{proof}

    \item[3] Suppose $𝑣, 𝑤 ∈ 𝑉$. Explain why there exists a unique $𝑥 ∈ 𝑉$ such that $𝑣 + 3𝑥 = w$.
    
    \begin{align*}
        v + 3x &= w \\
        3x &= w - v \\ 
        x &= \frac{1}{3}(w - v)
    \end{align*}
    Since that right side is a unique vector, $x$ is a unique vector. You invoke the properties and stuff for this, but I'm too lazy.

    \item[5] Show that in the definition of a vector space (1.20), the additive inverse
condition can be replaced with the condition that
$$0𝑣 = 0 \text{ for all } 𝑣 ∈ 𝑉.$$
Here the 0 on the left side is the number 0, and the 0 on the right side is the
additive identity of $V$.

    \begin{proof}
        We can show that given this condition, every element $v \in V$ has an additive inverse. 

        \begin{align*}
            0v &= 0 \\
            (1 + -1)v &= 0
        \end{align*}
        Note that even though we don't know whether \emph{vector} addition has an additive inverse the 0 on the left size is just a scalar, so we can still say there is an additive inverse. 


        Then distribute and apply the multiplicative identity axiom\dots
        \begin{align*}
            1v + (-1)v &= 0 \\
            v + (-1)v &= 0
        \end{align*}
        Therefore $(-1)v$ is the additive inverse for all $v \in V$.
    \end{proof}
    
    \item[7] Suppose $𝑆$ is a nonempty set. Let $𝑉^𝑆$ denote the set of functions from $𝑆$ to $𝑉$.
Define a natural addition and scalar multiplication on $𝑉^𝑆$
, and show that $𝑉^𝑆$
is a vector space with these definitions.

\begin{proof}
    Like with $F^S$, we can sort of ``piggyback'' off the properties of $V$ to maintain the properties of a vector space in $V^S$. 

    Define addition as 
    \[(f + g)(x) = f(x) + g(x)\]

    Let the additive identity be $0(x) = 0$ for all $x \in S$.

    \begin{align*}
        (f + 0)(x) = f(x) + 0(x) = f(x) + 0 = f(x)
    \end{align*}


    For all $f, g \in V^S$. Because $f(x)$ is a vector, it has an inverse (namely $-f(x)$). Therefore, we can define the additive inverse as 
    \[(-f)(x) = -f(x)\] 
    \begin{align*}
        (f + (-f))(x) = f(x) + (-f(x)) = 0 = 0(x)
    \end{align*}

    Communativity: 

    \begin{align*}
        (f + g)(x) = f(x) + g(x) = g(x) + f(x) = (g + f)(x)
    \end{align*}
    We can swap $f(x)$ and $g(x)$ because $V$ is communative.

    Associativity:
    \begin{align*}
        ((f + g) + h)(x) = (f(x) + g(x)) + h(x) = f(x) + (g(x) + h(x)) = (f + (g + h))(x)
    \end{align*}
    Again, we use the properties of the vector space $V$ to help us out here.

    Scalar multiplication 

    For all $\lambda \in \F, v \in V^S$, 
    \[(\lambda v)(x) = \lambda v(x)\]

    Multiplicative Identity 

    1 is the multiplicative identity. 
    \[(1v)(x) = 1v(x) = v(x)\]
    Since 1 is the multiplicative identity of $V$.

    Distributive property 

    \begin{align*}
        (a(u + v))(x) = a(u + v)(x) = a(u(x) + v(x)) = au(x) + av(x) = (au + av)(x)
    \end{align*}
    Because of the distributive property of $V$.

    \begin{align*}
        ((a + b)v)(x) = (a + b)v(x) = av(x) + bv(x) = (av + bv)(x)
    \end{align*}
    For the same reason.

    Thus, we have proved all the properties of vector spaces on $V^S$, therefore $V^S$ is a vector space.

    Closure 

    Because $V$ is closed under addition and scalar multiplication, and all our operations are defined using these 2 operations, $V^S$ is closed.

\end{proof}

\end{enumerate}

